%%%%%%%%%%%%%%%%%%%%%%%%%%%%%%%%%%%%%%%%%%%%%%%%%%%%%%%%%%%%%%%%%%%%%%%%%%%%%%%%
%%  INTRODUCTION
%%%%%%%%%%%%%%%%%%%%%%%%%%%%%%%%%%%%%%%%%%%%%%%%%%%%%%%%%%%%%%%%%%%%%%%%%%%%%%%%
\chapter{Introduction}\label{intro}
Since the introduction of Unimate in the early 1960s, robotic systems have been central to industrial manufacturing. In the automotive sector in particular, specialized robotic arms execute repetitive tasks such as painting, welding, and assembly with accuracy and speed. These operations, however, required a deterministic environment, so robots were confined to cages and pre-programmed paths, since any unplanned event could cause errors or safety hazards. Sensors and adaptive control changed this, enabling robots to absorb minor environmental variation \cite{Li2020}. The arrival of language models has pushed this further, as robots are beginning to navigate unfamiliar conditions with minimal human control, drawing on the world knowledge and reasoning these models provide.

This progress has enabled the field of \gls{eai}, where a model is embedded in a robotic body to close the loop between perception and action \cite{liu2025aligningcyberspacephysical}. Yet planning a task is not the same as executing it, and a gap remains between a model’s reasoning and a robot’s physical movement. Precise coordination is still difficult, and many models struggle with temporal consistency, failing to account for past physical states when planning future movements.

In this thesis, we bridge that disconnect with a framework that merges \gls{ros}~\cite{Macenski_2022} and MoveIt~\cite{coleman2014reducingbarrierentrycomplex} for motion planning, together with a \gls{yolo}-based~\cite{yolov8_ultralytics} vision pipeline for environmental awareness. To handle reasoning, we augment a \gls{vlm}~\cite{mistral_magistral_small_2025} with \gls{rag}~\cite{lewis2021retrievalaugmentedgenerationknowledgeintensivenlp} and an optional \gls{kg}~\cite{Hogan_2021}. Together, the stack can translate natural-language instructions into concrete, executable robot commands. Our novel contribution is an integrated architecture that combines planning, perception, and reasoning for coordinated robot control. The architecture aims at eventual \gls{s2r} deployment, but we build and validate it inside a Unity~\cite{unity} simulation. Ultimately, we aim to show that \glspl{vlm} can coordinate genuine, autonomous action between two robotic arms working toward a shared goal.

The rest of this chapter covers our motivation, objectives, and contributions. Chapter~\ref{background} summarizes related work, from early semantic planning to modern deployment frameworks. Chapter~\ref{methodology} lays out the system design and the reasoning behind it, from the operation system and the \gls{vlm}-driven command pipeline through to multi-robot coordination. Chapter~\ref{implementation} covers the implementation and the engineering decisions that mattered most. In Chapter~\ref{evaluation}, we define our benchmarks, present the results, and report an ablation for each component. Chapter~\ref{discussion} discusses the limitations of our evaluation and the open questions it leaves. Chapter~\ref{conclusion} closes with our contributions and future directions.

% ------------------------------------------------------------------------------
\section{Motivation}
Most \gls{eai} research focuses on single-agent tasks, yet many of the applications with the highest automation value are cooperative. Moving from a single arm to two introduces a new class of problems a lone robot never faces, because the arms must avoid colliding not only with objects but also with each other, and must synchronize their timing. Even with the rapid progress of language models, translating digital reasoning into reliable physical execution remains a challenge.

General-purpose models are also prone to hallucinating actions and lack the operational memory required by multi-stage task chains. To address this, we ground the model’s reasoning in live perception. A \gls{rag} operation pool restricts it to a vocabulary of executable actions, and a vision system resolves manipulation targets from current stereo detections, keeping its picture of the scene up to date.

This motivates a pipeline in which the reasoning model works as a high-level coordinator, selecting and ordering operations from a shared pool at runtime and translating intent into executable trajectories through \gls{ros} and MoveIt. Building this in a Unity-based simulation also creates a safe, scalable environment to develop these behaviors before they ever touch a physical motor.

% ------------------------------------------------------------------------------
\section{Goals, Novelty and Contributions}\label{sec:goals}\label{sec:novelty}
We set out to build an autonomous robotic operation framework driven by a reasoning \gls{vlm}, validated in simulation, and designed for physical deployment. The system accepts natural-language instructions and generates coordinated behavior without fine-tuning, hard-coded sequencing, or cloud-service dependencies. We build it in a Unity simulation with two AR4~\cite{AnninRobotics2026} arms, each behind a hardware abstraction layer.

The novelty is not a new model or control algorithm, as every component is based on previous work. The contribution is a runtime composition of behavior between two arms by a model that reasons over a shared vocabulary of basic operations, rather than a fixed execution path. The model selects, orders, and parameterizes the coordination primitives itself. The retrieval layer and the prompt supply handoff and synchronization priors, and the model applies them more reliably than it would if it had to rediscover them on its own (Section~\ref{sec:results-dual}). Prior systems side-step even that, assigning independent robots~\cite{ahn2024autortembodiedfoundationmodels} or negotiating through open dialogue~\cite{mandi2023rocodialecticmultirobotcollaboration}. Runtime composition of shared-object coordination from a fixed set of operations remains unexplored.

Five contributions build out this claim. First, a retrieval-grounded pipeline, built on a hierarchical registry of 20 operations, separates reasoning from execution. It combines \gls{rag} retrieval, a two-phase chain-of-thought decomposition inspired by RT-H~\cite{belkhale2024rthactionhierarchiesusing}, and a variable-passing mechanism that threads perception outputs into later operation parameters at runtime. Second, the coordination behind that claim runs on general-purpose primitives. The \gls{vlm} assembles signal/wait and parallel-group operations into time-ordered handoffs, parallel scheduling, and re-detection. The prompt supplies the handoff and synchronization rules, while the model chooses which primitives to use and how to combine them. Third, a consensus-gated negotiation protocol, inspired by RoCo~\cite{mandi2023rocodialecticmultirobotcollaboration} but without its free-form dialectic. Each robot gets its own agent, and the protocol runs a three-phase analysis–proposal–evaluation cycle. One agent drafts a complete plan per round in round-robin order; every peer must accept it, and the system falls back to single-model parsing if no consensus is reached within three rounds. Its B13 ablation shows a measurable but model-specific effect. For the default backend, per-task success rises $+10$~percentage points (pp), though the benefit does not generalize across models. Fourth, an operation-gated reflection mechanism re-queries the parser with an operation’s own recovery suggestions, for eligible operation categories only, clearing transient failures that a naive retry fixes (B12, $+14$~pp at the task level, $+7.0$~pp at the per-operation level). Fifth, an AutoRT adaptation for fixed-base manipulation swaps the raw-image scene representation for a structured object manifest and extends the prompt-encoded Robot Constitution with a hard kinematic layer (workspace bounds, velocity, gripper force, and minimum robot separation) that the semantic filter cannot override. Throughout, we isolate each component using single-flag ablations rather than reporting only aggregate success.

% ------------------------------------------------------------------------------
\section{Extensibility and Reusability}\label{sec:extensibility}
Because this system is intended to outlive the present thesis and serve as a foundation for later student projects, several engineering choices are designed to keep it extensible. The Python backend is split into strict dependency layers with centralized lazy imports, which avoids the circular-import failures that a growing research codebase accumulates and makes clear where new code belongs. New robot behavior has a single extension point. A contributor writes one \texttt{BasicOperation} with its parameter schema and optional metadata, registers it, and it propagates automatically to every subsystem that reads the registry as its single source of truth (Appendix~\ref{appendix:operation_example}). Camera and hardware interfaces sit behind adapters that isolate the rest of the system from the simulation/hardware distinction, so a move to physical hardware is localized to those adapters (Section~\ref{sec:sim-to-real}); tunable values live in dedicated configuration modules and Unity ScriptableObject \texttt{.asset} files (Appendix~\ref{appendix:config}); and each major subsystem is independently toggleable through its own feature flag. A regression test suite guards both halves of the system, with Python tests covering the operation registry, orchestrators, protocol handling, and the integration paths, while a Unity C\# test suite exercises the \gls{ik} solver, grasp pipeline, robot and trajectory controllers, and the wire protocol. Together, they give future contributors a safety net that catches breakage when they extend or refactor the system. The system also degrades gracefully when a piece is missing. Retrieval falls back to \gls{tfidf}, parsing to a regex matcher, negotiation to single-model parsing, and \gls{ros}/MoveIt motion to the built-in Unity controller, so the rest keeps running while one component is under development.